% !TeX root = ../papers/nc-p3.tex
\section{埋込み点を含む点曲線の余長と高次数の純性}
\label{sec:pure-curve-finite-defect}
点ファイバーを持つ標準有限Frobenius窓で共通点グラフ $X$ が一次元とする。
$\mathcal E\subset\mathcal O_X$ を有限台の最大イデアル、$Y$ をその商の純曲線とし、
$d$ を片側次数、$g=p_a(Y)$ とする。孤立したArtin環も $\mathcal E$ の長さに含む。
\begin{theorem}\label{thm:pure-curve-defect}
重複・可約成分や埋込み点を許して
\[
                    \ell\mathcal E=18-4d+2g .
\]
$d=7$ では $g=5,\ \mathcal E=0$、$d=6$ では $g=3,\ \mathcal E=0$ である。
六次の各辺で一般の五つの関係を取ると、五式内の残余曲線は両射影で型 $(2,2)$ の四次完全交叉となる。
$d=5$ なら $(g,\ell\mathcal E)=(1,0),(2,2)$ のいずれかであり、
$d=4$ なら $g=-1,0,1$ に限られる。
\end{theorem}
\begin{proof}
定理\ref{thm:point-graph-curve-degree}の一般五式の完全交叉 $W$ と、$Y$ の残余 $Z$ を取る。
$\omega_W=\mathcal O_W(1,1)$、$I_{Y/W}(1,1)=\omega_Z$ である。
第六式の切断 $\mathcal O_Z\to\omega_Z$ は各一般点で全射である。
実際そこで $X$ と $Y$ は等しく、$W$ 上の $I_X$ は第六式が生成する。
$\omega_Z$ と $\mathcal O_Z$ は各一般点で長さが同じなので、この射はそこでは同型である。
$Z$ は純ゆえ全体で単射となる。その余核を $(-1,-1)$ でひねったものは
$I_{Y/W}/I_{X/W}=\mathcal E$ である。従って長さは $2p_a(Z)-2$。
$\chi(\mathcal O_W)=-10$、$Z$ の両側次数 $10-d$ と上のイデアルのEuler標数から
$p_a(Z)=g+10-2d$、主張の式を得る。

純な次数4の非平面空間曲線 $V$ は、非被約でも算術種数1以下であり、
種数1なら二つ以上の独立な二次式があることを示す。一般平面断面の四点スキーム $A$ への
斉次座標環の商の全射を使う。長さ $n$ の平面零次元スキームのHilbert関数は
$n$ に達するまで各次数で少なくとも1増加する。一般一次非零因子で割った標準次数代数の
Hilbert関数がその一階差で、一度零なら以後も零になるためである。
元の座標環を一般平面式で割った環の次数1は3、次数2は3以上、次数3以降は4以上である。
もし次数2が3なら元の二次式は三本以上ある。しかし純四次を含む独立二次式が共通因子を持たなければ
型 $(2,2)$ で既に全次数を占め、三本目はない。いずれの二本も共通因子を持つ場合は
三次元の二次式空間に共通の一次因子がある。一本を固定してその高々二つの一次因子で割れる
部分空間の和集合を使えばよい。三本が $\ell a_j$ なら $\ell\mathcal O_V$ は三独立一次式の
共通点に支持され、純性により零となり平面に入ってしまう。
よって次数2は4以上、総和から $g(V)\le1$、等号なら元の次数2の次元は8である。
特に算術種数3の純四次は平面に入り、この性質は同じ次数の他の埋込みにも及ぶ。

$d=6$ の四次残余は先の単射で種数1以上である。平面なら両側で平面四次だが、
五式の基点は曲線 $W$ であり、両平面の積で一般の三式が切る完全交叉の両次数は3である。
従って平面でなく種数1、$\mathcal E=0$、$\omega_Z=\mathcal O_Z$ である。
二つの二次式が共通一次因子 $\ell$ を持てば、$(\ell)$ はその残りの二一次式の共通直線 $T$ 上の
$\mathcal O_T(-1)$ の商で、非平面性と純性からこれに等しい。
この単射は $\mathcal Hom_Z(\mathcal O_T,\omega_Z)=\omega_T=\mathcal O_T(-2)$ を通るので不可能である。
従って両側で完全交叉となる。七次の主張は定理\ref{thm:septic-point-graph}による。

五次の場合は $g\ge1$ である。$Y$ に純な平面四次部分曲線は含まれない。
その部分曲線の算術種数3と先の上界、成分ごとの偏極次数一定性から両位置で平面四次となり、
平面の積での切断に反する。$Y$ 自身が平面五次となる場合も除く。
実際、非平面なら上と同じ平面断面の議論で、一般平面式で割った環の次数 $0,1,2,3,4,\ldots$ の
次元が少なくとも $1,3,3,4,5,\ldots$ なので算術種数は5以下で、平面五次の種数6にならない。
従って平面なら両位置で平面でありやはり不可能となる。

非平面五次で $g\ge3$ と仮定する。二次式がなければ上の次数2の下界が6となり、
総和から $g\le2$ で矛盾する。既約な二次曲面に入れば滑らかな場合は有効二次数の和5から種数は高々2、
錐なら純な因子の類は母線と二超平面の和で、前節の七次と同じHilbert多項式計算から種数2となる。
ここで純な余次元一のイデアルは深さの完全列により正規曲面上で反射的である。
残る二次式を $\ell m$ とし、二因子の順を選ぶと $(\ell)$ の次数 $e$ は2以下となる。
実際 $(\ell)$ は $\mathcal O_{Y\cap(m=0)}(-1)$ の商であり、$(\ell),(m)$ の次数の和は5以下である。
非平面性から $e=1,2$ で、$(\ell)=\mathcal O_{D}(-1)$ の $D$ は平面 $m=0$ の純な次数 $e$ の因子である。
反対の商 $Y\cap(\ell=0)$ の平面内の純部分は次数 $5-e$ の因子で、他の有限余長を $b\ge0$ とする。
完全列から
\[
             g=p_a(D)+p_a(\text{平面次数 }5-e)+e-1-b .
\]
$e=2$ なら右辺は2以下、$e=1$ で $g\ge3$ なら純な平面四次部分曲線を持ってしまう。
これで五次を得る。四次の純部分自体も同様に平面ではなく、上界1と余長の非負性から結論する。
\end{proof}

\begin{proposition}\label{prop:embedded-fixed-length}
$Y$ が被約な局所完全交叉曲線で、四移動がどの成分上でも恒等でないとする。
$Z_{\rm fix}$ を $Y$ 上の固定スキームとすると
\[
        \ell Z_{\rm fix}=2p_a(Y)-2+\ell (\mathcal E/\mathcal E^2)|_Y .
\]
右の有限層は孤立部では零である。この式は固定スキームがCartierでなくても成立する。
\end{proposition}
\begin{proof}
$P=\mathbb P^3\times\mathbb P^3$、$H=L_0L_1$、
$J=I_Y/(I_X+I_Y^2)=\mathcal E/\mathcal E^2$ と置く。
$R_0 H^{-1}\to I_Y/I_Y^2$ は階数6から5への写像で、余核が $J$ である。
その核を $S H^{-1}$ と書く。余法束は被約曲線の一般点で周囲の微分に単射なので
$S$ は微分の飽和核でもある。定理\ref{thm:fixed-differential}の反対評価切断は
$S L_2 L_{-1}$ に入り、その係数は $Y$ 上の固定イデアル $\mathfrak f$ を生成する。
実際後ろの辺を $Y\times Y$ に制限したファイバー積は $Y$ のグラフそのものである。
一般点でこの切断が核を張るので、飽和核は
$S=(L_2 L_{-1})^{-1}\mathcal Hom(\mathfrak f,\mathcal O_Y)$ である。
局所的には切断の有理倍率で六係数全てを正則にするものを取ったことになる。
Gorenstein曲線上の零次元層の局所双対により、$\mathcal Hom(\mathfrak f,\mathcal O_Y)$ は
$\mathcal O_Y$ を余長 $\ell Z_{\rm fix}$ で含む。余法束の行列式は $H^{-4}\omega_Y^{-1}$。
上の四項完全列の次数（無捩れ層ではEuler標数による次数）を取り、
$\deg L_2 L_{-1}=\deg H$ を代入すれば表示式となる。
\end{proof}

\begin{corollary}\label{cor:integral-hull-high-degrees}
曲線に重なる有限構造も許し、純部分 $Y$ が整な被約曲線であるとする。
標準AS窓で $d\ge5$ なら $X=Y$ で、楕円五次の二類と八係数の整七次だけが現れる。
$d=4$ なら $Y$ は種数1の $(2,2)$ 完全交叉で、節点・尖点や埋込み点を許しても
定理\ref{thm:finite-orbit-quartic-marks}が全次数の必要十分判定を与える。
\end{corollary}
\begin{proof}
六次・七次の純性と整曲線の分類を使う。五次では算術種数1なら純性が従う。
種数2のGorenstein曲線では定理\ref{thm:singular-integral-quintic-exclusion}の
二次曲面の残余母線から得る三次鉛筆の議論を適用する。
一般の滑らかな $p$ を取る二核束、$Q_2$ への非零制限、$D_p+r_p$ の計算は
有限台の追加構造を避け、全て $Y$ 上の切断と完全線型系のみによる。
従って四移動は二重被覆の対合で固定長は6である。前命題では
$\ell J=6-2=4$ を必要とするが、$J$ は長さ2の $\mathcal E$ の商なので不可能である。
非Gorensteinなら同定理の唯一の錐頂点は $X$ 上でも接次元3である
（$Y$ で既に3、$X$ は三空間に埋まる）。自動Ore表示の水平射影と極次数上限は
他の孤立点・埋込み点があっても同じである。曲面はなく、可動な水平二次曲線の上の高さは
一般点で一意なので非Gorenstein五次も消える。

整四次の算術種数は0または1である。種数0なら $Y=\mathbb P^1$ で、
$\nu^*\Omega_X$ の捩れ長は常に $\ell J$ に一致する。
実際 $Y$ が滑らかなら余法列は局所分裂し、微分の余核は $\Omega_Y$ に加えて $J$ を持つ。
非恒等な射影直線自己同型の固定長は2、恒等の場合は定理\ref{thm:fixed-differential}の第二項から、
どちらでも $\ell J=4$ となる。ところが $\ell\mathcal E=2$。
残る整な種数1の曲線は滑らかまたは唯一の節点・尖点を持ち双対化束は自明である。
ゆえに記載の四次印の定理を適用できる。
\end{proof}
